\paragraph{Preparing the Input State}\label{ex:preparing-the-input-state}

As we have seen in the previous section, Deutsch's algorithm uses two qubits:
\begin{itemize}
    \item The \textbf{data qubit}, which represents the input $x$;
    \item The \textbf{computation qubit}, which allows the oracle to encode $f(x)$ reversibly.
\end{itemize}
The algorithm begins with:
\begin{equation*}
    \ket{\psi_0} = \ket{0}\ket{1}
\end{equation*}
The first qubit starts in the state \ket{0} (data qubit), while the second starts in \ket{1} (computation qubit).

\highspace
\textcolor{Red2}{\faIcon{exclamation-triangle}} At this point, we could apply the oracle $U_f$ to the state \ket{\psi_0}. However, this would not be useful. Indeed:
\begin{equation*}
    U_f \ket{\psi_0} = U_f \ket{0, 1} = \ket{0, 1 \oplus f(0)}
\end{equation*}
And if $f(0) = 1$:
\begin{equation*}
    U_f \ket{\psi_0} = \ket{0, 1 \oplus 1} = \ket{0, 0}
\end{equation*}
The second qubit now literally stores the answer, but the Deutsch's algorithm is \textbf{not interested in the value of $f(0)$}, because it \hl{does not tell us whether the function is constant or balanced!}

\highspace
\textcolor{Green3}{\faIcon{check-circle}} So we need to prepare the input state in a way that allows the oracle to encode the function's behavior on both inputs $x=0$ and $x=1$. This is achieved by putting the data qubit in a \textbf{superposition} of \ket{0} and \ket{1}. Then a Hadamard gate is applied to both qubits:
\begin{equation*}
    \ket{\psi_1} = \left(H \otimes H\right) \ket{0, 1}
\end{equation*}
Using the definition of the Hadamard gate:
\begin{equation*}
    H\ket{0} = \ket{+} = \frac{\ket{0} + \ket{1}}{\sqrt{2}}, \qquad H\ket{1} = \ket{-} = \frac{\ket{0} - \ket{1}}{\sqrt{2}}
\end{equation*}
Where \ket{+} represents the \emph{data qubit} and \ket{-} represents the \emph{computation qubit}. We can rewrite the state \ket{\psi_1} as:
\begin{align*}
    \ket{\psi_1} &= \left(H \otimes H\right) \ket{0, 1} \\
    &= H\ket{0} \otimes H\ket{1} \\
    &= \ket{+}\ket{-} \\
    &= \frac{\ket{0} + \ket{1}}{\sqrt{2}} \otimes \frac{\ket{0} - \ket{1}}{\sqrt{2}} \\
    &= \frac{1}{2} \left(\ket{0, 0} - \ket{0, 1} + \ket{1, 0} - \ket{1, 1}\right)
\end{align*}
This preparation is not arbitrary. Each of the two states has a precise purpose.

\newpage

\begin{flushleft}
    \textcolor{Green3}{\faIcon{question-circle} \textbf{Why is the data qubit prepared in $\ket{+}$?}}
\end{flushleft}
The data qubit must \hl{allow the oracle to act on both possible inputs}:
\begin{equation*}
    x = 0 \quad \text{and} \quad x = 1
\end{equation*}
Within the same quantum evolution (i.e., in superposition). The state:
\begin{equation*}
    \ket{+} = \frac{\ket{0} + \ket{1}}{\sqrt{2}}
\end{equation*}
Is an \textbf{equal superposition of the two inputs}. Therefore, when the oracle is applied, linearity gives us:
\begin{align*}
    U_f \left(\ket{+} \otimes \ket{-}\right) &= U_f \ket{+} \otimes U_f \left(\ket{-}\right) \\
    &= U_f \left(\frac{\ket{0} + \ket{1}}{\sqrt{2}}\right) \otimes U_f \left(\ket{-}\right) \\
    &= \left[\frac{1}{\sqrt{2}} U_f \left(\ket{0} + \ket{1}\right)\right] \otimes U_f \left(\ket{-}\right) \\
    &= \frac{1}{\sqrt{2}} U_f \left(\ket{0, -} + \ket{1, -}\right)
\end{align*}
We simply rewrote the state to demonstrate that the same application of the oracle acts on both branches of the superposition: \ket{0} and \ket{1}. In other words, it only means: ``\emph{dear oracle, we want to ask about \textbf{both possible inputs} $x=0$ and $x=1$, so evaluate the function for 0 and 1 simultaneously}''. Thus, the firs qubit is \textbf{not the asnwer}; it is simply a \textbf{superposition of the two possible questions} we want to ask the oracle.

\highspace
However, we should be careful with the common statement that the quantum computer ``\emph{evaluates both values simultaneously}''. The \textbf{algorithm does not directly reveal both} $f(0)$ and $f(1)$. Measurement cannot extract both outputs. Instead, the \textbf{oracle encodes a relation} between them \textbf{into the relative phase} of the data qubit.

\highspace
\textcolor{Green3}{\faIcon{question-circle} \textbf{Why is the data qubit necessary if it doesn't do anything?}} At first sight, the data qubit may appear to play a passive role because the oracle never changes its computational basis state. Indeed, in:
\begin{equation*}
    U_f \ket{x,y} = \ket{x, y \oplus f(x)}
\end{equation*}
The value of $x$ is preserved. However, this does \textbf{not} mean that the data qubit remains unchanged. What changes are the \textbf{relative phases of its amplitudes}.

\highspace
Before the oracle, the data qubit is prepared in the equal superposition:
\begin{equation*}
    \ket{+} = \frac{\ket{0} + \ket{1}}{\sqrt{2}}
\end{equation*}
After the oracle, it becomes:
\begin{equation*}
    \frac{(-1)^{f(0)}\ket{0} + (-1)^{f(1)}\ket{1}}{\sqrt{2}}
\end{equation*}
Notice that the basis states are still \ket{0} and \ket{1}; the oracle has \textbf{not} changed the value of $x$. Instead, it has modified the \textbf{relative phase} between the two branches of the superposition. This relative phase is where the information about the function is encoded.

\highspace
\textcolor{Green3}{\faIcon{question-circle} \textbf{But how can the oracle modify the data qubit without directly acting on it?}} The answer lies in the \textbf{computation qubit}. Because it is prepared in the eigenstate \ket{-}, the oracle converts the classical output bit $f(x)$ into a phase factor $(-1)^{f(x)}$, which is then transferred to the corresponding branch of the data qubit through \textbf{phase kickback}. Let's see how this works in detail.

\highspace
\begin{flushleft}
    \textcolor{Green3}{\faIcon{question-circle} \textbf{Why is the computation qubit prepared in $\ket{-}$?}}
\end{flushleft}
Recall that the oracle acts as:
\begin{equation*}
    U_f \ket{x, y} = \ket{x, y \oplus f(x)}
\end{equation*}
For firxed $x$, the oracle therefore either (\autopageref{ex:oracle-implementations-for-the-four-functions}):
\begin{itemize}
    \item Does nothing to the second qubit when $f(x) = 0$;
    \item Applies $X$ to the second qubit when $f(x) = 1$.
\end{itemize}
Indeed:
\begin{equation*}
    U_f \ket{x,y} = \ket{x} \otimes X^{f(x)} \ket{y}
\end{equation*}
Where $X^{f(x)}$ is the identity when $f(x) = 0$ and the $X$ gate when $f(x) = 1$:
\begin{equation*}
    X^{f(x)} = \begin{cases}
        X^{0} = I & \text{if } f(x) = 0 \\
        X^{1} = X & \text{if } f(x) = 1
    \end{cases}
\end{equation*}
We want the computation qubit to react differently to these two cases, but without storing the result as an ordinary bit. This is achieved by choosing an eigenstate of $X$, because eigenstates are only multiplied by a phase factor when acted upon by their corresponding operator. In other words, we want the computation qubit to \hl{convert the classical output bit $f(x)$ into a detectable relative sign}.

\highspace
The Pauli-$X$ eigenstates are:
\begin{equation*}
    X\ket{+} = +\ket{+}, \quad \text{and} \quad X\ket{-} = -\ket{-}
\end{equation*}
The state \ket{-} has eigenvalue $-1$. Therefore:
\begin{equation*}
    X^{f(x)} \ket{-} = \begin{cases}
        X^{0} \ket{-} = I\ket{-} = \ket{-} & \text{if } f(x) = 0 \\
        X^{1} \ket{-} = X\ket{-} = -\ket{-} & \text{if } f(x) = 1
    \end{cases}
\end{equation*}
This can be written more compactly as:
\begin{equation*}
    X^{f(x)} \ket{-} = (-1)^{f(x)} \ket{-}
\end{equation*}
Since the data qubit is in the superposition:
\begin{equation*}
    \ket{+}
    =
    \frac{\ket{0}+\ket{1}}{\sqrt{2}}
\end{equation*}
Linearity allows us to write the action of the oracle on the two-qubit state as:
\begin{equation*}
    \begin{aligned}
        U_f(\ket+\ket-)
        &=
        \frac1{\sqrt2}
        \left(
        U_f\ket{0,-}
        +
        U_f\ket{1,-}
        \right)\\
        &=
        \frac1{\sqrt2}
        \left(
        (-1)^{f(0)}\ket{0,-}
        +
        (-1)^{f(1)}\ket{1,-}
        \right)\\
        &=
        \left(
        \frac{
        (-1)^{f(0)}\ket0
        +
        (-1)^{f(1)}\ket1
        }
        {\sqrt2}
        \right)
        \otimes
        \ket-
    \end{aligned}
\end{equation*}
This expression \hl{explains why Deutsch's algorithm never learns the individual values $f(0)$ and $f(1)$.} The \textbf{oracle only modifies the relative sign between the two branches}:
\begin{equation*}
    \frac{\ket0+\ket1}{\sqrt2}
    \quad\longrightarrow\quad
    \frac{
    (-1)^{f(0)}\ket0
    +
    (-1)^{f(1)}\ket1
    }{\sqrt2}
\end{equation*}
Only the relative phase matters physically.
\begin{itemize}
    \item \important{Constant case}. If the two function \textbf{values} are \textbf{equal}, both amplitudes receive the \textbf{same sign}, which corresponds to the state $\ket{+}$ (up to an irrelevant global phase):
    \begin{equation*}
        \begin{aligned}
            f(0) = f(1) = 0 &\implies \frac{\ket0+\ket1}{\sqrt2} = \ket{+} \\
            f(0) = f(1) = 1 &\implies \frac{-\ket0-\ket1}{\sqrt2} = -\frac{\ket0+\ket1}{\sqrt2} = -\ket{+} = \ket{+}
        \end{aligned}
    \end{equation*}
    Notice that the global phase $-1$ is physically irrelevant, so we can ignore it.
    \item \important{Balanced case}. If the two function \textbf{values} are \textbf{different}, the amplitudes acquire \textbf{opposite signs}, which corresponds to the state $\ket{-}$ (again up to a global phase):
    \begin{equation*}
        \begin{aligned}
            f(0) = 0, f(1) = 1 &\implies \frac{\ket0-\ket1}{\sqrt2} = \ket{-} \\
            f(0) = 1, f(1) = 0 &\implies \frac{-\ket0+\ket1}{\sqrt2} = -\frac{\ket0-\ket1}{\sqrt2} = -\ket{-} = \ket{-}
        \end{aligned}
    \end{equation*}
    Again, the global phase $-1$ is physically irrelevant, so we can ignore it.
\end{itemize}
Consequently, the \textbf{oracle's action} can be summarized as:
\begin{equation}
    U_f \ket{x}\ket{-} = (-1)^{f(x)} \ket{x}\ket{-}
\end{equation}
Where:
\begin{itemize}
    \item $\ket{x}$ is the data qubit, which explores both inputs $x=0$ and $x=1$ simultaneously;
    \item $\ket{-}$ is the computation qubit, which converts the classical output $f(x)$ into a relative phase factor $(-1)^{f(x)}$ (phase kickback). Its role is therefore to \hl{convert $f(x) \in \left\{0,1\right\}$ into the phase factor $(-1)^{f(x)} \in \left\{+1,-1\right\}$}.
\end{itemize}
\textcolor{Green3}{\faIcon[regular]{lightbulb}} Note that the computation qubit is in the state \ket{-} because \hl{its negative eigenvalue transforms the classical output bit $f(x)$ into a detectable relative sign.}

\highspace
In summary, the two qubits are prepared differently because they perform two different jobs:
\begin{itemize}
    \item \ket{+} explores both inputs $x=0$ and $x=1$ simultaneously;
    \item \ket{-} converts the classical output $f(x)$ into a relative phase factor $(-1)^{f(x)}$ (phase kickback).
\end{itemize}